% =====================================================================
\section{АЦП: читаем аналоговые сигналы}\label{les:adc}
% =====================================================================
\lessonintro
  {научиться измерять напряжение, понять ограничения АЦП ESP32-S3, измерять напряжения выше 3,3 В}
  {урок 2 (вывод результатов в Serial), урок 4 (яркость через ШИМ)}
  {потенциометр управляет яркостью светодиода}
  {потенциометр \SI{10}{k\ohm}; для раздела про делитель --- резисторы \SI{100}{k\ohm} и \SI{33}{k\ohm}}

\subsection{Сначала: сигнал не всегда только 0 или 1}

Потенциометр делит 3,3 В между GND и GPIO: в середине получается около 1,65 В. АЦП измеряет это
напряжение и возвращает число. Так программа понимает, насколько повернули ручку. \code{int} хранит
целое число; \code{analogReadMilliVolts} читает милливольты; 1000 мВ = 1 В.

\subsection{Как работает АЦП}

АЦП последовательного приближения (SAR) преобразует напряжение в число. При разрешении 12 бит
диапазон разбивается на $2^{12} = 4096$ уровней:
\[
  N = \left\lfloor \frac{U_\text{вх}}{U_\text{полн}} \cdot 4095 \right\rfloor .
\]

\begin{figure}[H]
\centering
\begin{tikzpicture}
\begin{axis}[
  width=0.62\textwidth, height=5.6cm,
  xlabel={$U_\text{вх}$, В}, ylabel={код АЦП},
  xmin=0, xmax=3.3, ymin=0, ymax=4300,
  grid=major, grid style={gray!25},
  legend pos=south east, legend style={font=\scriptsize},
  tick label style={font=\scriptsize}, label style={font=\small},
  scaled y ticks=false, yticklabel style={/pgf/number format/1000 sep={}},
]
  \addplot[very thick,gray,dashed,domain=0:3.3,samples=34] {min(x/3.1*4095,4095)};
  \addlegendentry{идеальный}
  \addplot[very thick,esp,smooth] coordinates {
    (0,0) (0.05,0) (0.1,40) (0.3,330) (0.8,1000) (1.5,1930) (2.2,2850) (2.7,3470) (2.95,3790) (3.1,3990) (3.2,4095) (3.3,4095)};
  \addlegendentry{реальный (без калибровки)}
  \draw[<-,thin] (axis cs:0.08,20) -- (axis cs:0.5,1400) node[above,font=\scriptsize,align=center]{«мёртвая зона»\\ у нуля};
  \draw[<-,thin] (axis cs:3.2,4095) -- (axis cs:2.4,3900) node[left,font=\scriptsize]{насыщение};
\end{axis}
\end{tikzpicture}
\caption{Характеристика АЦП (иллюстрация). Края диапазона нелинейны --- используйте калиброванное чтение}
\end{figure}

\subsubsection{Особенности АЦП ESP32-S3}

\begin{itemize}
  \item Два блока: \textbf{ADC1} (\pin{1}--\pin{10}) и \textbf{ADC2} (\pin{11}--\pin{20}).
  \item \textbf{ADC2 используется также Wi-Fi.} Доступ зависит от режима и драйвера. В уроке 5 плата будет подключена к сети,
        поэтому потенциометр ставим на ADC1 --- \pin{1}.
  \item Внутреннее опорное напряжение $\approx\SI{1.1}{V}$, поэтому на входе стоит \emph{аттенюатор}.
        По умолчанию в Arduino выбран \SI{11}{dB} --- рабочий диапазон примерно \SIrange{0}{3.1}{V}.
  \item Каждый чип откалиброван на заводе (eFuse). Функция \code{analogReadMilliVolts()} учитывает калибровку
        и возвращает милливольты --- предпочитайте её «сырому» \code{analogRead()}.
\end{itemize}

\subsection{Потенциометр --- «крутилка»}

\textbf{Потенциометр} --- это резистор с третьим, подвижным контактом (движком). Крайние выводы подключают
к 0 и 3,3~В, а движок, скользя по резистивной дорожке, выдаёт любое напряжение между ними. Поворот
ручки --- самый наглядный аналоговый сигнал.

\begin{twoview}{Потенциометр на выводе GPIO1}{fig:pot}
\begin{tikzpicture}[bb]
  \bbBoard{24}
  \bbDevKit{29}{25}
  \bbPotD{8}{12}{10k}
  \draw[wire=wblack] (8,14) -- (8,16);
  \draw[wire=wred] (10,14) -- (10,17);
  \draw[wire=wgreen] (dk-1) -- (39.5,22) -- (39.5,28.6) -- (9,28.6) -- (9,14);
  \draw[wire=wred] (dk-3V3) -- (27.5,25) -- (27.5,17) -- (24,17);
  \draw[wire=wblack] (dk-GND) -- (26.5,4) -- (26.5,16) -- (24,16);
\end{tikzpicture}
\nextview
\begin{circuitikz}
  \draw (0,3) node[vcc]{3,3 В} to[pR, n=pot, l_=\SI{10}{k\ohm}] (0,0) node[ground]{};
  \draw (pot.wiper) -- ++(1,0) to[short,-o] ++(0.6,0) node[right]{\pin{1} (ADC1)};
\end{circuitikz}
\end{twoview}

\codecap{Чтение потенциометра}
\codeinput{l11_pot}

\begin{tip}
Открой \textsf{Tools → Serial Plotter}: пары «имя:значение» Arduino IDE нарисует графиками. Покрути ручку ---
увидишь свой аналоговый сигнал вживую!
\end{tip}

\subsection{Проверка чтения ручки}

До подключения LED откройте Serial Monitor и медленно поверните ручку. Ожидается: около 0 мВ у GND,
около 3100 мВ у 3V3 (дальше АЦП упирается в предел диапазона) и промежуточные числа между ними. Если всегда 0, проверьте провод с ползунка
потенциометра до ADC. Если всегда максимум, проверьте, что крайние выводы действительно идут к GND и 3V3.

\begin{tip}
Откройте \textsf{Tools → Serial Plotter}: пары «имя:значение» IDE нарисует графиками. Покрутите
ручку и обратите внимание на шум --- соседние отсчёты отличаются на единицы-десятки милливольт.
\end{tip}

\subsection{Самостоятельный опыт: ручка яркости}

Переведите милливольты в проценты яркости и поборите шум двумя приёмами:
\begin{itemize}
  \item \textbf{усреднение} --- 8 измерений подряд;
  \item \textbf{порог} --- яркость меняется, только если ручку повернули больше чем на 3 \%,
        чтобы LED не дрожал при неподвижной ручке.
\end{itemize}

Свяжите уровень потенциометра с \code{ledcWrite()} из урока 4 в новом самостоятельном скетче.

\subsection{Измерение напряжения выше 3,3 В}

Чтобы измерить, например, батарею \SI{12}{V}, понижаем напряжение делителем:
\[
  U_\text{АЦП} = U_\text{вх}\cdot\frac{R_2}{R_1+R_2}
  = 12 \cdot \frac{33}{100+33} \approx \SI{2.98}{V}.
\]

\begin{figure}[H]
\centering
\begin{circuitikz}
  \draw (0,4) node[left]{$U_\text{вх}$ до 12 В} to[short,o-] (1,4)
        to[R, l=$R_1$, a=\SI{100}{k\ohm}] (1,2) coordinate(m)
        to[R, l=$R_2$, a=\SI{33}{k\ohm}] (1,0) node[ground]{};
  \draw (m) -- (3,2) coordinate(n) to[C, l=\SI{100}{nF}] (3,0) node[ground]{};
  \draw (n) to[short,-o] (4.5,2) node[right]{\pin{2}};
\end{circuitikz}
\caption{Делитель напряжения с фильтрующим конденсатором}
\end{figure}

\codecap{Вольтметр с усреднением}
\codeinput{l5_voltmeter}

\subsection{Практика: ручка поворачивает сервопривод}

Соединим «крутилку» и сервопривод: повернул ручку --- повернулся моторчик. Так устроены многие
манипуляторы роботов.

У сервопривода три провода: \textcolor{wbrown}{\textbf{коричневый}} --- земля, \textcolor{wred}{\textbf{красный}} ---
питание 5~В, \textcolor{worange}{\textbf{оранжевый}} --- сигнал. Питание сервопривода возьмём с вывода 5V
платы, поэтому на макетной плате у нас будут \emph{две разные шины «+»}: сверху 3,3~В для потенциометра,
снизу 5~В для мотора. Будь внимателен и не перепутай их! Шины «−» соединены перемычкой слева.

\begin{twoview}{Потенциометр управляет сервоприводом}{fig:servo}
\begin{tikzpicture}[bb]
  \bbBoard{24}
  \bbDevKit{29}{25}
  \bbPotD{8}{12}{10k}
  \draw[wire=wblack] (8,14) -- (8,16);
  \draw[wire=wred] (10,14) -- (10,17);
  \draw[wire=wgreen] (dk-1) -- (39.5,22) -- (39.5,28.6) -- (9,28.6) -- (9,14);
  \draw[wire=wred] (dk-3V3) -- (27.5,25) -- (27.5,17) -- (24,17);
  \draw[wire=wblack] (dk-GND) -- (26.5,4) -- (26.5,16) -- (24,16);
  \draw[wire=wblack] (1,16) -- (1,1);
  \draw[wire=worange] (dk-5V) -- (28.1,5) -- (28.1,0) -- (24,0);
  \draw[wire=wyellow] (dk-6) -- (27.7,19) -- (27.7,6) -- (18,6);
  \bbServoBody{12}{-4.5}
  \draw[line width=1.3pt,wbrown,-{Circle[length=3pt,fill=wbrown!70!black]}] (14,-4.5) -- (14,1);
  \draw[line width=1.3pt,wred,-{Circle[length=3pt,fill=wred!70!black]}] (15,-4.5) -- (15,0);
  \draw[line width=1.3pt,worange,-{Circle[length=3pt,fill=worange!70!black]}] (16,-4.5) -- (16,-2.5) -- (18,-2) -- (18,3);
\end{tikzpicture}
\nextview
\begin{circuitikz}
  \draw (0,3) node[vcc]{3,3 В} to[pR, n=pot, l_=\SI{10}{k\ohm}] (0,0) node[ground]{};
  \draw (pot.wiper) -- ++(1,0) to[short,-o] ++(0.4,0) node[right]{\pin{1}};
  \node[draw,thick,fill=blkper,minimum width=2.2cm,minimum height=1.8cm,font=\sffamily] (s) at (10,1.5) {Серво};
  \draw[worange,thick] ($(s.west)+(0,0.5)$) -- ++(-3,0) node[left,text=black]{\pin{6}};
  \draw[wred,thick] ($(s.west)+(0,0)$) -- ++(-3,0) node[left,text=black]{5V};
  \draw[wbrown,thick] ($(s.west)+(0,-0.5)$) -- ++(-1,0) -- ++(0,-0.6) node[ground]{};
\end{circuitikz}
\end{twoview}

\codecap{Ручка поворачивает сервопривод}
\codeinput{l11_pot_servo}

Функция \code{map(x, 0, 3100, 0, 180)} пересчитывает число из одного диапазона в другой: 0~мВ → 0°,
3100~мВ → 180°. А проверка \code{abs(angle - lastAngle) >= 2} не даёт мотору дёргаться из-за шума.

\begin{caution}
Под нагрузкой сервопривод потребляет сотни миллиампер. Если при его движении плата перезагружается,
питай сервопривод от отдельного источника 5~В (например, 4 батарейки АА), соединив его «минус» с GND платы.
\end{caution}


\subsection{То же на ESPUI: вольтметр со стрелкой и графиком}

Число в Serial Monitor удобно для отладки, но смотреть на ручку потенциометра и на график
одновременно приятнее на экране телефона. Если у вас уже собрана схема этого урока, дополнительных
деталей не нужно: соберите панель по образцу урока 5 (скетч ниже) и подключитесь к сети \code{S3-Voltmeter}
(пароль \code{lesson-39-demo}).

\codecap{ESPUI: Gauge и Graph для потенциометра на GPIO1}
\codeinput{l39_espui_voltmeter}

Крутите ручку: стрелка Gauge показывает милливольты, а на графике остаётся след последних
отсчётов. Усреднение по 16 замерам из этого урока работает прежним образом, а страница
получает новое значение два раза в секунду.

\begin{trouble}
\begin{itemize}[leftmargin=*]
  \item \textbf{Всегда 0 или всегда максимум} --- ползунок потенциометра не дошёл до GPIO или крайние выводы подключены неверно.
  \item \textbf{Показания сильно скачут} --- усредняйте 8--32 отсчёта и поставьте конденсатор 100 нФ на вход.
  \item \textbf{С Wi-Fi чтение перестало работать} --- датчик стоит на ADC2; перенесите его на ADC1 (\pin{1}--\pin{10}).
\end{itemize}
\end{trouble}

\begin{summary}
\begin{itemize}[leftmargin=*]
  \item \code{analogReadMilliVolts()} даёт калиброванный результат; диапазон --- примерно \SIrange{0}{3.1}{V}.
  \item Для совместной работы с Wi-Fi выбираем ADC1; для сетевых проектов берите ADC1 (\pin{1}--\pin{10}).
  \item Шум убираем усреднением, а дрожание управления --- порогом (гистерезисом).
\end{itemize}
\end{summary}

\begin{selfcheck}
\begin{enumerate}
  \item Почему \code{analogReadMilliVolts()} лучше, чем \code{analogRead()}?
  \item Почему в сетевых опытах аналоговый датчик подключают к ADC1?
  \item Как измерить напряжение 12 В и что делает конденсатор в делителе?
\end{enumerate}
\end{selfcheck}

\begin{tasks}
\begin{enumerate}
\item \lv{1} Выведите «сырое» значение АЦП потенциометра и убедитесь, что оно меняется от 0 до 4095.
\item \lv{1} Выведите напряжение в милливольтах через \code{analogReadMilliVolts()} и сравните с мультиметром в трёх положениях ручки.
\item \lv{2} Добавьте усреднение по 8 отсчётам. Измерьте, на сколько уменьшился разброс показаний.
\item \lv{2} Управляйте яркостью светодиода (урок 4) потенциометром и сделайте «мёртвую зону» в крайних положениях.
\item \lv{3} Добавьте фоторезистор в делитель и режим «авто»: светодиод ярче, когда в комнате темнее, с гистерезисом.
\item \lv{4} Измерьте \code{micros()}, сколько длится чтение с усреднением по 8 и по 64 отсчётам, и постройте график «точность против скорости».
\end{enumerate}
\end{tasks}
